\documentclass{article}
\usepackage[backend=biber,natbib=true,style=alphabetic,maxbibnames=50]{biblatex}
\addbibresource{/home/nqbh/reference/bib.bib}
\usepackage[utf8]{vietnam}
\usepackage{tocloft}
\renewcommand{\cftsecleader}{\cftdotfill{\cftdotsep}}
\usepackage[colorlinks=true,linkcolor=blue,urlcolor=red,citecolor=magenta]{hyperref}
\usepackage{amsmath,amssymb,amsthm,enumitem,float,graphicx,mathtools,tikz}
\usetikzlibrary{angles,calc,intersections,matrix,patterns,quotes,shadings}
\allowdisplaybreaks
\newtheorem{assumption}{Assumption}
\newtheorem{baitoan}{Bài toán}
\newtheorem{cauhoi}{Câu hỏi}
\newtheorem{conjecture}{Conjecture}
\newtheorem{corollary}{Corollary}
\newtheorem{dangtoan}{Dạng toán}
\newtheorem{definition}{Definition}
\newtheorem{dinhly}{Định lý}
\newtheorem{dinhnghia}{Định nghĩa}
\newtheorem{example}{Example}
\newtheorem{ghichu}{Ghi chú}
\newtheorem{hequa}{Hệ quả}
\newtheorem{hypothesis}{Hypothesis}
\newtheorem{lemma}{Lemma}
\newtheorem{luuy}{Lưu ý}
\newtheorem{nhanxet}{Nhận xét}
\newtheorem{notation}{Notation}
\newtheorem{note}{Note}
\newtheorem{principle}{Principle}
\newtheorem{problem}{Problem}
\newtheorem{proposition}{Proposition}
\newtheorem{question}{Question}
\newtheorem{remark}{Remark}
\newtheorem{theorem}{Theorem}
\newtheorem{vidu}{Ví dụ}
\usepackage[left=1cm,right=1cm,top=5mm,bottom=5mm,footskip=4mm]{geometry}
\def\labelitemii{$\circ$}
\DeclareRobustCommand{\divby}{%
	\mathrel{\vbox{\baselineskip.65ex\lineskiplimit0pt\hbox{.}\hbox{.}\hbox{.}}}%
}
\setlist[itemize]{leftmargin=*}
\setlist[enumerate]{leftmargin=*}

\title{Learn, Teach, \& Research Rules \& Principles\\Vài Quy Tắc \& Nguyên Lý về Học, Dạy, \& Nghiên Cứu}
\author{Nguyễn Quản Bá Hồng\footnote{A scientist- {\it\&} creative artist wannabe, a mathematics {\it\&} computer science lecturer of Department of Artificial Intelligence {\it\&} Data Science (AIDS), School of Technology (SOT), UMT Trường Đại học Quản lý {\it\&} Công nghệ TP.HCM, Hồ Chí Minh City, Việt Nam.\\E-mail: {\sf nguyenquanbahong@gmail.com} {\it\&} {\sf hong.nguyenquanba@umt.edu.vn}. Website: \url{https://nqbh.github.io/}. GitHub: \url{https://github.com/NQBH}.}}
\date{\today}

\begin{document}
\maketitle
\begin{abstract}
	This text is a part of the series {\it Some Topics in Advanced STEM \& Beyond}:
	
	{\sc url}: \url{https://nqbh.github.io/advanced_STEM/}.
	
	Latest version:
	\begin{itemize}
		\item {\it Learn, Teach, \& Research Rules \& Principles -- Vài Quy Tắc \& Nguyên Lý về Học, Dạy, \& Nghiên Cứu}.
		
		PDF: {\sc url}: \url{https://github.com/NQBH/advanced_STEM_beyond/blob/main/teach/rule_principle/NQBH_learn_teach_research_rule.pdf}.
		
		\TeX: {\sc url}: \url{https://github.com/NQBH/advanced_STEM_beyond/blob/main/teach/rule_principle/NQBH_learn_teach_research_rule.tex}.
	\end{itemize}
\end{abstract}
\tableofcontents

%------------------------------------------------------------------------------%

\section{Learn}
Bài tập phải có 1 danh sách đặc tả tất cả các yêu cầu của giảng viên, như khi đi làm phải có 1 danh sách đặc tả tất cả các yêu cầu của khách hàng vậy.

\paragraph{Requirements.}
\begin{enumerate}
	\item Dịch đề sang Tiếng Việt: đọc hiểu sâu \& giải thích đề (must have).
	\item Giải/Chứng minh (có thể dùng AI để hỗ trợ, nhưng phải kiểm tra lại nhiều lần để đảm bảo kết quả đúng hoặc tốt nhất có thể): check kỹ.
	\item Minh họa: Dùng hình ảnh, biểu đồ, văn xuôi để minh họa ví dụ trực quan (có thể sử dụng ẩn/hoán dụ nếu muốn, miễn là giải thích bằng cách hiểu của chính bản thân mình \& đúng là được).
	\item Nếu đã đạt 3 yêu cầu trên. Có thể mở rộng, tổng quát bài toán, vẫn phải tuân thủ 3 bước trên để nhận thêm nhiều điểm cộng.
	\item Không bắt làm hết, làm nhiều nhất có thể \& phải đảm bảo đúng/tốt nhất có thể. Viết report bằng LaTeX: max 10/problem, Word: max 8/problem, nếu xung phong lên bảng trong buổi học đó, không giấy + điện thoại thì max 15, trình bày viết bảng . Điểm tùy theo độ khó của bài toán \& mức độ hiểu của SV.
	\item Nộp mail: {\tt hong.nguyenquanba@umt.edu.vn}. Chỉ nộp 1 lần duy nhất. Mỗi lần nộp lại $-1$. Đặt sai tên file $-1$/file.
	\item Cách đặt tên file: tên viết tắt in hoa của SV + gạch chân + BT + số của BT + {\tt.{tex,pdf,doc,cpp,c,py}}. Sai filename - 1 điểm/file.
	\item Nén tất cả file bài tập lại thành 1 file với cách đặt tên file: NNTA.rar hoặc NNTA.zip xong nộp về Mail:\\{\tt hong.nguyenquanba@umt.edu.vn} bằng official mail sinh viên.
	\item Chỉ được nộp 1 lần duy nhất, nên check cẩn thận trước khi nộp, mỗi lần nộp lại -1.
	\item Phải chú thích từng dòng code.
	\item Những file code nào generated by AI >= \verb|AI_percentage_maximum_threshold|\% sẽ bị trừ điểm rất nặng với giá trị mặc định \verb|AI_percentage_maximum_threshold| $= 20$.
	\item Nộp giống code chia đôi điểm nếu điểm dương, trừ như nhau nếu điểm âm.
	\item Chửi thề: $-10$.
	\item Chửi hoặc có hành vi vô giáo dục với giảng viên: full 0.
\end{enumerate}

%------------------------------------------------------------------------------%

\section{Teach}

%------------------------------------------------------------------------------%

\section{Research}

%------------------------------------------------------------------------------%

\section{Miscellaneous}

%------------------------------------------------------------------------------%

\printbibliography[heading=bibintoc]
	
\end{document}